\section{方法}
\label{sec:method}
策略只能看到 2,048 行训练数据和它自己调用的结果。它维护当前最好点 $W^\ast$（目前实测最好的点；初始为最佳训练行，该行不会被
重新发送），每次调用提出 50 个新点。

\subsection{加性分块下降}
\label{sec:block}
记 $B_t$ 为第 $t$ 次调用的\emph{基点}，$B_t[b\!\leftarrow\!w]$ 为把基点的第 $b$ 块替换为 $w$ 后的点。除下文的组合点外，每个候选
都是这种单块形式，因此其实测效果 $\delta_b(w)=R(B_t[b\!\leftarrow\!w])-R(B_t)$ 只属于一栋建筑。每栋建筑维护一张表 $v_b$，记录
实测过的块及其相对当前最好点的效果。在加性假设下，
\begin{equation}
\hat R(w_1,w_2,w_3)=R(W^\ast)+\textstyle\sum_{b=1}^{3}\bigl(v_b(w_b)-v_b(w^\ast_b)\bigr),
\label{eq:add}
\end{equation}
它不需要拟合任何模型。下一次调用的基点，是在当前最好点的块和每栋建筑两个最好的改进块所构成的组合中 $\hat R$ 最低的那个；
另有最多六个组合与它一起发送。这一设计同时带来三点好处：每个点都是可归因的实验；一个批次并行推进三个 7 维搜索；不同建筑
中找到的改进只需各花一个点就能组合起来。它是一种分块坐标下降~\cite{bcd}，像加性贝叶斯优化~\cite{addgp}那样利用加性，但用的
是实测的、而非模型推断的块效果。

在一个块内，候选来自四类\emph{移动}：\emph{代理模型}挑选，从 1,500 个扰动中取下置信界（均值 $-$ 标准差）最低者，代理模型由
岭回归趋势加上基于距基点最近 600 行的 ARD GP 组成；\emph{随机}移动，采用逐建筑的信任尺度，成功后乘以 1.5、失败后乘以 0.6，
限定在 $[0.01,1]$ 内~\cite{turbo}；\emph{精细}移动（小的乘性步长、把一个次要权重置零，或设为训练权重 0.1\,\% 与 10\,\% 分位数
之间的极小份额，$1.6\times10^{-6}$ 到 $7.2\times10^{-3}$）；以及从第 9 次调用起的\emph{扫描}，即在已知平台之外一点的单坐标步，
外加把上一次调用中有改进的坐标重新施加到新基点上。第 1--6 次调用中，偏移移动偏向那些在训练行中份额与 $R$ 负相关的目标
（变体 v2）。

\begin{algorithm}[t]
\caption{策略的一次调用（v4）}
\label{alg:call}
\small
\begin{algorithmic}[1]
\Require 训练行、本次运行已测得的点、种子、调用序号 $t$
\State $W^\ast\gets$ 实测最好点；更新块表 $v_b$、信任尺度、零效应区间、命中计数 $h_c,n_c$
\State $B_t\gets$ 在每栋建筑两个最好块的组合中取 $\hat R$ 最小者~\eqref{eq:add}
\State 批次 $\gets\{B_t\}\,\cup$ 最多六个次优组合
\State 按~\eqref{eq:share} 计算各类名额份额 $s_c$
\For{建筑 $b=1,2,3$ 与类别 $c$}
\State 从 $B_t[b]$ 出发抽取 $s_c n_b$ 个第 $c$ 类的块 $w$（$n_b$：建筑 $b$ 的剩余名额，取整），跳过零效应区间和已发送点
\State 批次 $\gets$ 批次 $\cup\,\{B_t[b\!\leftarrow\!w]\}$
\EndFor
\State 发送 50 个点，记录每个 $R$ \Comment{随机数以（种子, $t$）为种子}
\end{algorithmic}
\end{algorithm}

\subsection{证据驱动的分配}
\label{sec:alloc}
固定的类别份额会把大部分名额花在收获很少的地方：份额为 50/25/25 时，代理模型挑选占了 44\,\% 的点，命中率却只有 3.5\,\%。
变体 v4 增加了三条只使用本次运行结果的规则。\emph{平台过滤}：一个单坐标移动若恰好返回 $R(B_t)$（$|\delta|\le10^{-9}$），就为该
建筑、基点块和坐标标记一个零效应区间；此后不再在最宽的此类区间内提出移动。\emph{后验份额}：若最近四次调用中第 $c$ 类的
$n_c$ 个点里有 $h_c$ 个命中（低于其调用的基点），则
\begin{equation}
s_c=f_c+\Bigl(1-\textstyle\sum_{c'}f_{c'}\Bigr)\frac{r_c}{\sum_{c'}r_{c'}},\qquad r_c=\frac{h_c+1}{n_c+5},
\label{eq:share}
\end{equation}
即 Beta(1,4) 后验均值，代理模型挑选的下限 $f_c=0.25$，其他类为 0.10，使任何一类都不会被饿死。\emph{不重复}：与训练行或已发送
点相同的候选从不提出，所以每次调用都发送 50 个新点。

\subsection{部署}
\label{sec:deploy}
提交的策略（h7）在第 1--11 次调用中就是 v4；从第 12 次调用起，它把最多 12 个精细或置零名额替换为更粗的单坐标、多坐标和重启
移动。所有随机性都由（种子，调用）派生；评分运行包固定了 BLAS 内核和线程数，因为 GP 超参数拟合会把舍入差异放大成不同的
选择。运行包在发送每个批次前先把它写入磁盘，绝不重发已扣额度的调用，并且不调用仿真器就能从日志重新计算每一个批次。
